%
% 52-kettenintegral.tex -- Integral über Ketten
%
% (c) 2026 Prof Dr Andreas Müller
%
\subsection{Integral über Ketten}
Damit das Wegintegral eine lineare Abbildung von Ketten wird, 
muss das Wegintegral entlang einer Kette durch Linearkombination
der einzelnen Integrale gebildet werden:

\begin{definition}[Wegintegral einer Kette]
Ist $c\in C_1$ eine Kette der Form
\eqref{buch:integration:homologie:eqn:def:kette}
und $f$ eine in $U$ holomorphe Funktion, dann ist das 
\emph{Wegintegral entlang $C$} definiert durch
\index{Wegintegral entlang einer Kette}%
\index{Kettenintegral}%
\begin{equation}
\int_c f(z)\,dz
=
c_1
\int_{\gamma_1}f(z)\,dz
+
\ldots
+
c_n
\int_{\gamma_n}f(z)\,dz
=
\sum_{k=1}^n
c_k \int_{\gamma_k} f(z)\,dz.
\label{buch:integration:homologie:eqn:kettenintegral}
\end{equation}
\end{definition}

Ist $F(z)$ eine Stammfunktion von $f(z)$, dann kann das Integral über
die Komponenten $\gamma_k$ der Kette ausgerechnet werden und es
gilt
\[
\int_c f(z)\,dz
=
\sum_{k=1}^n c_k \int_{\gamma_k}f(z)\,dz
=
\sum_{k=1}^n c_k
\Bigl[F(z)\Bigr]_{\gamma_k(0)}^{\gamma_k(1)}
=
\sum_{k=1}^n c_k
\bigl(
F(\gamma_k(1)) - F(\gamma_k(0))
\bigr)
\]
Die Klammer im letzten Term kann man interpretieren als die
Auswertung der Funktion $F$ auf dem Rand des Wegstückes $\gamma_k$.
Schreibt man dies rein formal als das Integral
\[
\int_{\partial\gamma_k} F(z)\,dz
=
\Bigl[
F(z)
\Bigr]_{\gamma_k(0)}^{\gamma_k(1)},
\]
dann kann man die Gleichung
\eqref{buch:integration:homologie:eqn:kettenintegral}
als
\begin{equation}
\int_{\partial c} F(z)\,dz
=
\int_c F'(z)\,dz
\label{buch:integration:homologie:eqn:stokes0}
\end{equation}
schreiben.
Diese Formel ist ein Beispiel für den allgemeinen Satz von Stokes
\index{Satz!von Stokes}%
aus der Theorie der Differentialformen.
\index{Differentialform}
Er besagt, dass das Integral der Ableitung einer Differentialform
über ein Gebiet ersetzt werden kann durch das Integral der Differentialform
über den Rand.
In der Schreibweise der Differentialformen wird dies für eine
Differentialform $\omega$ und das Gebiet $S$ als
\[
\int_{\partial S}\omega
=
\int_S d\omega
\]
geschrieben, $d\omega$ ist die äussere Ableitung der Differentialform.
Die Ähnlichkeit mit
\eqref{buch:integration:homologie:eqn:stokes0}
ist unverkennbar.

%
% Integral über Zyklen
%
\subsubsection{Integral über Zyklen}
Ist $c$ ein Zyklus, dann lassen sich die Teilwege von $c$ zu einem
geschlossenen Weg zusammen setzen, wie dies in
Abbildung~\ref{buch:integration:homologie:fig:kettenzyklen}
dargestellt ist.
In diesem Fall ist es sinnvoll vom Integral über einen geschlossen
Weg zu sprechen und
\[
\partial c = 0
\qquad\Rightarrow\qquad
\int_c f(z)\,dz
=
\oint_c f(z)\,dz
\]
zu schreiben.

Abbildung~\ref{buch:integration:homologie:fig:paar}
zeigt eine Situation einer Kette $C=\gamma_++\gamma_-$, deren
Wegintegral
\[
\int_C f(z)\,dz
=
\int_{\gamma_+} f(z)\,dz
+
\int_{\gamma_-} f(z)\,dz
=
\int_{\gamma} f(z)\,dz
\]
mit dem Wegintegral entlang der Kurve $\gamma$ übereinstimmt.
Soweit es um Wegintegrale entlang einer Kette geht, gibt es 
zwischen $\gamma$ und $C$ keinen Unterschied.

Zwischen $\gamma$ und nur der Komponente $\gamma_+$ gibt es
ein wesentlichen Unterschied.
Die Funktion $f(z) = 1/(z-F_1)$ hat eine Polstelle im Punkt $F_1$.
Die Wegintegrale
\begin{align*}
\frac{1}{2\pi i}
\oint_\gamma f(z)\,dz
&=
\\
\frac{1}{2\pi i}
\oint_\gamma \frac{1}{z-F_1}\,dz
&=
\\
\operatorname{ind}_\gamma(F_1)
&=
1
\\
&\ne
\frac{1}{2\pi i}
\oint_{\gamma_+} f(z)\,dz
\\
&=
\frac{1}{2\pi i}
\oint_{\gamma_+} \frac{1}{z-F_1}\,dz
\\
&=
\operatorname{ind}_{\gamma_+}(F_1)
\\
&=
0
\end{align*}
sind jedoch verschieden.
Der Grund dafür ist, dass $\gamma$ und $\gamma_+$ den Punkt $F_1$
nicht gleich oft umlaufen.
Wir definieren daher ein Äquivalenzrelation von Ketten durch die
Anzahl der Male, die sie Punkte ausserhalb von $U$ umlaufen.

\begin{definition}[Umlaufzahl einer Kette]
Die \emph{Umlaufzahl} einer Kette $C$ um einen Punkt $z_0$ ist
\index{Umlaufzahl}%
\[
\operatorname{ind}_C(z_0)
=
\frac{1}{2\pi i}
\int_C
\frac{1}{z-z_0}
\,dz.
\]
\end{definition}
